\begin{abstract}
We characterize minimax squared-error precision for a population-weighted treatment contrast with binary treatment and outcomes and finite discrete covariates. The experiment supplies independent samples from a common population: \(n\ge1\) complete records containing outcomes, treatment, and covariates, and \(m\ge0\) auxiliary treatment--covariate records. Covariates have \(d\ge2\) categories, and treatment propensities in occupied cells lie in \([\epsilon,1-\epsilon]\), where \(\epsilon\in(0,1/4]\) is a public overlap floor. Over arbitrary covariate probabilities and cellwise outcome means, the minimax risk is
\[
\asymp
\min\left\{
1,\frac{1}{n\epsilon}
+\left[
\frac{d}{(n+m)\epsilon\log(\exp(1)+n\epsilon)}
\right]^2
\right\},
\]
with numerical comparison constants independent of all four indices. A deterministic estimator on the original records attains this order with proof-oriented tuning and numerical constants, and a matching lower bound applies to all Borel randomized rules on those records. The comparison yields exact sequence criteria for consistency, for joint consistency and rare-label benchmark attainment, and for a vanishing risk ratio relative to complete records alone. At fixed complete-record size, alphabet size, and overlap floor, increasing auxiliary information converges to the exact-marginal minimax benchmark. Under potential-outcome consistency and conditional treatment exchangeability, the contrast equals the population average treatment effect.
\end{abstract}

\section{Introduction}

This paper establishes a uniform minimax precision frontier for estimating a population-weighted treatment contrast when outcome information and treatment--covariate information arrive through separate sampling channels. The model has binary treatment and outcomes and a finite covariate alphabet. Its unknown population law permits arbitrary covariate probabilities and unrestricted cellwise outcome means, subject to a public treatment-overlap restriction in every occupied covariate cell. The resulting risk characterization describes jointly how the complete-record budget, the auxiliary-record budget, the alphabet size, and a possibly diminishing overlap floor determine attainable precision.

The experiment contains \(n\ge1\), the number of complete outcome--treatment--covariate records, and \(m\ge0\), the number of independent auxiliary treatment--covariate records, sampled from the same population. The covariate alphabet has size \(d\ge2\), and \(\epsilon\in(0,1/4]\), the public overlap floor, bounds occupied-cell treatment propensities away from zero and one. Thus all \(n+m\) records carry information about population treatment--covariate frequencies, while the complete sample also carries outcomes. Conditional overlap accommodates arbitrarily small population covariate cells, making the distinction between these information budgets consequential when many cells are sparsely represented.

For this experiment, \cref{obj:thm:uniform-annotation-frontier} establishes that \(R(n,m,d,\epsilon)\), the minimax squared-error risk, satisfies
\[
R(n,m,d,\epsilon)
\asymp
\min\left\{
1,\frac{1}{n\epsilon}
+\left[
\frac{d}{(n+m)\epsilon\log(\exp(1)+n\epsilon)}
\right]^2
\right\}.
\]
The comparison constants are numerical and independent of every public index throughout the stated domain. The target is the population-covariate-weighted difference between the two conditional outcome means, defined in \cref{obj:def:target}. Every observable model law admits a compatible potential-outcome extension satisfying consistency and conditional treatment exchangeability, and the target equals the population average treatment effect under every such extension, by \cref{obj:lem:causal-realization}.

The two terms identify distinct sources of statistical difficulty. The inverse scale \(1/(n\epsilon)\) reflects information about rare-arm outcomes in the complete sample. The many-cell term reflects learning the unknown treatment--covariate distribution using both channels. A common-marginal testing construction in \cref{obj:thm:rare-label-floor} establishes a capped inverse-\(n\epsilon\) risk floor for every auxiliary budget, including the experiment supplied with the exact population marginal. The deterministic construction in \cref{obj:def:estimator-handle} attains the full frontier through independent observation pools, pilot localization, reciprocal-polynomial correction, clipping, and finite averaging over ordered prefixes. Its proof-oriented tuning establishes the statistical order uniformly over the model and public indices with comparison constants that absorb the zero-output threshold; \cref{sec:attaining-estimator} explains the mechanism.

The frontier also gives exact resource characterizations for arbitrary experiment sequences. Writing \(S_k=n_k\epsilon_k\) for rare-arm outcome information, \(N_k=n_k+m_k\) for the total marginal-information budget, and \(\ell_k=\log(\exp(1)+S_k)\) for the regularized logarithmic scale, \cref{obj:thm:annotation-resource-phases} shows that consistency is equivalent to
\[
S_k\longrightarrow\infty,
\qquad
d_k=o(N_k\epsilon_k\ell_k).
\]
Joint consistency and attainment of rare-label benchmark order are equivalent to
\[
S_k\longrightarrow\infty,
\qquad
d_k=O\!\left(\frac{N_k\epsilon_k\ell_k}{\sqrt{S_k}}\right).
\]
The same result characterizes a vanishing risk ratio relative to complete records alone, accounting for both supervised saturation and supervised risks that already vanish. These criteria quantify the marginal expansion needed to reach a precision order governed by complete-record outcomes. At fixed \(n,d,\epsilon\), \cref{obj:prop:benchmark-recovery} further establishes convergence of the minimax risk itself to the exact-marginal risk as \(m\to\infty\), with an explicit finite-sample comparison.

The closest supervised analysis is \citet[model (3), Theorems 1--2]{ZengBalakrishnanHanKennedy2026Discrete}, and the matching fixed-interior supervised characterization is credited to \citet[Theorem \texttt{thm:sharp-minimax-fixed-interior}]{InternalDiscreteATEPolynomialUpperMatch}. The fixed-overlap two-channel frontier is established by \citet[Theorem 4]{CausalSmith2026Annotation}. Relative to these comparisons, the present result characterizes the joint dependence on unequal information budgets and a possibly diminishing public overlap floor through constants uniform over all four indices. Reciprocal approximation in off-policy evaluation provides a methodological antecedent \citep[Section 3.3, Theorems 5--6]{MaZhuJiaoWainwright2022OPE}. Semi-supervised treatment-effect inference supplies complementary efficiency results under nuisance-estimation conditions \citep[Theorem 2.1 and Corollary 2.1]{ChakraborttyDai2024Semisupervised}, while the allocation results of \citet[Theorems 4.2 and 5.1 of the arXiv manuscript]{NwankwoGoldkindZhou2026Annotations} address annotation design under their adaptive sampling regime. Here the precision comparison concerns passive independent acquisition from a common population.

The results hold under the stated model and sampling assumptions.

The paper proceeds from related work to the experiment, target, and causal interpretation, followed by the uniform minimax comparison and benchmark recovery. It then explains the attaining estimator and develops the consistency and auxiliary-information resource characterizations before discussing their interpretation and limitations. The appendices present the causal and benchmark arguments, estimator risk bounds, testing and approximation lower bounds, resource derivations and verification note, and proofs of the main results.
